\section{多指标证据与采样过程分析}
\subsection{相同标量行预算下的独立误差}
为避免只根据相对解误差评价算法，表\ref{tab:multimetric}补列1232行主实验的最大绝对误差、两侧归一化PDE残差、界面值RMS和物理通量RMS。全部为五种子的算术均值；原始逐次数据与标准差仍保存于同一CSV及汇总表。归一化PDE残差除以各自扩散系数，而最后一列通量未作归一化，读表时不可把这些量直接按数值大小比较。
\begin{table}[p]\centering
\caption{相同1232行预算下的独立多指标均值；每侧64特征、scale2、三轮}\label{tab:multimetric}
\resizebox{\textwidth}{!}{\begin{tabular}{lllllll}\toprule
$(\beta_1,\beta_2)$ & 方法 & $E_\infty$ & $R_1/\beta_1$ & $R_2/\beta_2$ & $J_u$ & $J_q$ \\ \midrule
(1,1) & A & $4.895\times10^{-4}$ & $2.989\times10^{-3}$ & $3.125\times10^{-3}$ & $1.124\times10^{-4}$ & $9.465\times10^{-5}$ \\
(1,1) & B & $4.449\times10^{-4}$ & $3.866\times10^{-3}$ & $4.382\times10^{-3}$ & $1.185\times10^{-4}$ & $1.179\times10^{-4}$ \\
(1,1) & C & $5.999\times10^{-4}$ & $3.227\times10^{-3}$ & $3.010\times10^{-3}$ & $1.519\times10^{-4}$ & $1.236\times10^{-4}$ \\
(1,1) & D & $6.116\times10^{-4}$ & $2.974\times10^{-3}$ & $2.988\times10^{-3}$ & $1.452\times10^{-4}$ & $1.074\times10^{-4}$ \\
(1000,1) & A & $4.537\times10^{-4}$ & $3.108\times10^{-3}$ & $2.656\times10^{-3}$ & $1.123\times10^{-4}$ & $1.029\times10^{-1}$ \\
(1000,1) & B & $3.964\times10^{-4}$ & $4.613\times10^{-3}$ & $3.906\times10^{-3}$ & $1.093\times10^{-4}$ & $1.047\times10^{-1}$ \\
(1000,1) & C & $5.385\times10^{-4}$ & $3.075\times10^{-3}$ & $2.574\times10^{-3}$ & $1.450\times10^{-4}$ & $1.376\times10^{-1}$ \\
(1000,1) & D & $5.506\times10^{-4}$ & $3.036\times10^{-3}$ & $2.502\times10^{-3}$ & $1.436\times10^{-4}$ & $1.219\times10^{-1}$ \\
(1,1000) & A & $4.082\times10^{-4}$ & $2.665\times10^{-3}$ & $2.913\times10^{-3}$ & $7.398\times10^{-5}$ & $8.728\times10^{-2}$ \\
(1,1000) & B & $3.720\times10^{-4}$ & $3.402\times10^{-3}$ & $3.868\times10^{-3}$ & $7.790\times10^{-5}$ & $9.671\times10^{-2}$ \\
(1,1000) & C & $4.923\times10^{-4}$ & $2.733\times10^{-3}$ & $2.961\times10^{-3}$ & $9.955\times10^{-5}$ & $1.135\times10^{-1}$ \\
(1,1000) & D & $5.153\times10^{-4}$ & $2.655\times10^{-3}$ & $2.869\times10^{-3}$ & $9.029\times10^{-5}$ & $1.095\times10^{-1}$ \\
\bottomrule\end{tabular}
}
\end{table}

在$\beta=(1000,1)$设置中，D的右侧归一化PDE残差均值约$2.50\times10^{-3}$，低于B的约$3.91\times10^{-3}$；但D的最大绝对误差约$5.51\times10^{-4}$，高于B的约$3.96\times10^{-4}$。这提供了一个具体反例：在独立点上内部残差降低，也不能不加条件地推出最大解误差降低。原因可能同时涉及界面、边界、截断及不同误差传播方向，不能只根据这两列唯一归因。

同一设置下，D的物理通量RMS约0.122，B约0.105，说明残差目标在内部改善时并未同时改善全部界面指标。这个量级来自每侧64特征的采样主比较，不能与每侧128特征的界面验证表混用。若不注明特征数，读者可能误以为高对比度通量结果前后矛盾；实际上两表属于不同近似能力和配置点实验。

对于单位系数，D两侧PDE残差也处于较低水平，但界面值误差和最大解误差没有同步胜过静态基线。因此本研究的负面结果不仅是某个高对比度下的偶然异常。另一方面，不能把当前三个平滑设置概括为联合采样在所有问题上无效；峰值实验已经显示D相对仅PDE采样在部分设置改善，这种差异仍需按问题和指标分别报告。

\subsection{条件数较低为何仍可能误差较高}
表\ref{tab:rankcost}显示相同行预算的条件数、有效秩、物理点及时间。这里总列数为128，所有列出的主比较有效秩均为128，原始条件数约为$10^8$量级。D在高对比度设置下的平均条件数略低于B，却仍没有更小的解误差。这说明条件数反映最坏线性敏感性，不能直接作为近似误差的单调代理。
\begin{table}[p]\centering
\caption{1232行主比较的矩阵与资源诊断，五种子；时间含完整自适应过程}\label{tab:rankcost}
\resizebox{\textwidth}{!}{\begin{tabular}{lllllll}\toprule
$(\beta_1,\beta_2)$ & 方法 & $\kappa$均值 & 有效秩 & 物理点均值 & 界面点均值 & 时间/秒 \\ \midrule
(1,1) & A & $1.812\times10^{8}$ & 128--128 & 1104.0 & 128.0 & 0.0135 \\
(1,1) & B & $1.766\times10^{8}$ & 128--128 & 1104.0 & 128.0 & 0.0137 \\
(1,1) & C & $1.618\times10^{8}$ & 128--128 & 1104.0 & 128.0 & 0.0752 \\
(1,1) & D & $1.636\times10^{8}$ & 128--128 & 1168.4 & 63.6 & 0.0755 \\
(1000,1) & A & $1.314\times10^{8}$ & 128--128 & 1104.0 & 128.0 & 0.0167 \\
(1000,1) & B & $1.311\times10^{8}$ & 128--128 & 1104.0 & 128.0 & 0.0132 \\
(1000,1) & C & $1.238\times10^{8}$ & 128--128 & 1104.0 & 128.0 & 0.0773 \\
(1000,1) & D & $1.183\times10^{8}$ & 128--128 & 1168.6 & 63.4 & 0.0837 \\
(1,1000) & A & $1.507\times10^{8}$ & 128--128 & 1104.0 & 128.0 & 0.0139 \\
(1,1000) & B & $1.516\times10^{8}$ & 128--128 & 1104.0 & 128.0 & 0.0142 \\
(1,1000) & C & $1.365\times10^{8}$ & 128--128 & 1104.0 & 128.0 & 0.0752 \\
(1,1000) & D & $1.348\times10^{8}$ & 128--128 & 1169.2 & 62.8 & 0.0759 \\
\bottomrule\end{tabular}
}
\end{table}

首先，矩阵条件数不包含真解到特征空间的距离。一个条件良好但空间表达能力不足的系统，也可能给出准确求解的错误近似。其次，相同行预算不保证相同的约束经验分布；矩阵列的可辨识性略有改善，可能伴随对界面约束覆盖的改变。再次，真解误差对应特定右端及预测位置，最坏方向未必被显著激发。

这与第二章固定矩阵误差分解一致：除了扰动放大项，还存在验证近似、截断和求值项。当前主比较未截断出明显秩亏，不能据此排除空间与经验目标的影响；后续增宽实验出现的有效秩不足则属于另一种情形。统一使用条件数解释所有实验，会混淆不同宽度和截断状态。

时间方面，自适应方案约为多次求解及候选评估之和，而A/B只求解一次。表中小规模CPU时间受系统状态影响，不宜将单次差值解释为精确硬件常数；但完整成本显著增加且解误差没有改善，足以限制当前方案的实用性主张。更严格的固定峰值串行重测已单独报告，不能拿并发阶段的时间去宣称加速。

\subsection{配额轨迹提供的机制证据}
表\ref{tab:history}列出$\beta=(1000,1)$时D的三轮平均轨迹，均值不必是整数。每次实际运行中的配额均为整数并满足预算，不能把均值中的小数理解为程序使用了小数个点。
\begin{table}[htbp]\centering
\caption{联合策略D的实际平均配额轨迹；每侧64特征、五种子}\label{tab:history}
\begin{tabular}{lllll}\toprule
预算 & 轮次 & 左内部均值 & 右内部均值 & 界面均值 \\ \midrule
点 & 0 & 128.0 & 128.0 & 32.0 \\
点 & 1 & 212.4 & 225.6 & 42.0 \\
点 & 2 & 307.0 & 310.2 & 54.8 \\
点 & 3 & 399.0 & 397.0 & 68.0 \\
行 & 0 & 128.0 & 128.0 & 32.0 \\
行 & 1 & 223.0 & 238.6 & 41.2 \\
行 & 2 & 331.2 & 332.8 & 52.0 \\
行 & 3 & 438.6 & 426.6 & 63.4 \\
\bottomrule\end{tabular}

\end{table}

点预算下最终界面平均68点，少于静态128点，而内部各侧平均接近400点；行预算下界面平均约63.4点。对应物理点均值超过1104，是因为一个界面位置占两行，减少界面后可用同样行预算放更多内部位置。这一计数关系解释了表\ref{tab:rankcost}中D的物理点更多，却没有使用额外标量行。

轨迹支持“当前指标倾向把资源给内部”这一描述，但不单独证明界面点少就是误差变大的唯一原因。改变配额同时改变位置分布、组经验平均和拟合结果；还可能存在某些内部点重复提供近似信息。后续同配额固定随机对照只回答一项有限问题，不能完成全部机制识别。

如果以后改进分配规则，首先应规定如何平衡物理通量与归一化通量，以及是否按每组单位成本的收益分配。当前按残差幅值分配并未测量“再增加一点能使真解误差降低多少”，两者不是同一量。由于真实应用没有真解，任何收益估计也必须来自可计算的代理，而不能把本实验已知解析解偷偷用于选点。

\section{参数敏感性结果的进一步解释}
宽度、权重范围和点数不是可以独立保证收敛的三个旋钮。增宽扩展名义特征空间，但若新增列高度相关，可辨识维数增长有限；增加点数可能改善覆盖，却因组平均权重改变和空间误差重新分配而导致解误差非单调；改变范围既改变特征尺度，也改变导数和矩阵谱。因此单因素扫描回答的是固定其余设置时的局部响应，不是完整的联合最优化。

在平滑直线模型上较小范围表现较好，可以解释为某些低频结构容易由缓变特征组合表示；同时小范围也可能增强列相关。这两种现象并不冲突。若只报告误差，可能忽略系数病态；若只报告条件数，又可能错过一个在指定双精度环境下确实准确的近似。本文把二者并列，避免用单个数字替代数值分析。

界面平方权重增加可加强满足界面约束的倾向，但可能牺牲内部拟合。它与增加界面点数并不等价：前者直接改变组目标的重要性，后者在组均方归一化下主要改变经验覆盖。把二者都笼统称为“加强界面损失”，会掩盖具体机制。本文所有对照均保留原权重与点数，便于单独核对。

岭参数也不是越大越稳定越好。较强收缩可以减小系数，但可能去掉真解所需方向；未经截断和经过岭的条件数必须区分数据矩阵与增广矩阵。现有扫描没有建立自动选择岭参数的方法，也没有用最终解析误差为每个种子挑参数，因此报告的是固定设置的敏感性，而非可在未知真解下直接部署的参数选择规则。

\section{近似复现中的可重复与可对齐}
可重复是指用本项目公开配置能够再次得到相近输出；与原文可对齐则要求使用相同模型、初始化、点分布、导数、权重、求解阈值和误差规则。前者并不自动推出后者。本文已经保存自身配置，但缺少作者若干实现细节，因此保留近似复现定位。

原表误差与本项目差异不能通过任意更换种子、删去较差运行或只保留较小积分值来消除。先前对初始化范围的另一解释也保留为诊断，而没有改写成作者原始方案。这种处理可能使结果表更复杂，却能使读者知道哪些设置来自原文字面，哪些来自本项目的补充假设。

解析导数与FD导数在部分设置接近，只说明在当前初始化、阈值及误差量级下，两种实现得到相近预测；不能证明FD矩阵与解析矩阵完全相同，也不能推导所有范围和步长下等价。步长扫描显示很小的差分步长可能恶化误差，符合舍入放大的机制，但其定量幅度还受线性系统敏感性影响。

花形积分误差的相对变化在最小误差组较大，提示接近极小数值时应降低对末几位的解释强度。本文不比较两方法某一末位数字来宣称优劣，也不把数值积分结果当作严格误差认证。后续若需高精度结论，应增加积分密度和高精度求值，并明确把额外成本计入。

\section{对原研究设想的回答}
原设想是利用内部与界面残差共同分配资源，以便在界面困难或高对比度时更有效地采样。现有实现确实完成了这一机制，也完成了相应基线及消融；因此研究过程仍围绕原问题展开。变化在于结果没有支持预期的稳定优势，论文据此把贡献从候选新方法收束为独立实现与系统分析。

已有证据支持三个有限判断：一是合理静态配额有时已能达到消融收益，二是内部残差下降不保证解误差和物理通量同步下降，三是自适应本身有不可忽略的完整成本。这些判断都能在表格中找到对应配置，不依赖作者主观挑选最好运行。它们也不等价于证明更精细的联合指标永远无用。

后续最有价值的改进应针对已观察的机制提出可证伪假设，例如在固定表达空间下分离配额与位置选择，或在冻结权重后使用新的监测池。若改成增长网络、更多子域或更大特征数，必须给静态基线相同新增表达能力，否则无法知道收益是否来自采样。先明确比较对象，再决定是否值得继续扩展，比不断扩大参数扫描更有解释力。

本章最终没有以“算法一定成功”作为验收标准，而是以模型与代码一致、资源计数清楚、全种子数据保留和结论可核对作为标准。尚未满足的学校形式要求及文献缺口属于终稿准备事项，不能通过把负面结果改成正面陈述来解决。
